\documentclass[12pt,a4paper]{article}

\usepackage{graphicx}
\usepackage{tikz}
\usepackage{hyperref}
\usepackage{epstopdf}
\usepackage{amsmath}
\usepackage{pdfpages}
\usepackage{indentfirst}
\usepackage{url}
\usepackage{float}
\usepackage{multirow}
\usepackage{booktabs}
\usepackage{caption}

%\newcommand{\sym}[1]{\rlap{$#1$}} % for symbols in Table
\newcommand{\sym}[1]{{#1}} % for symbols in Table
%\graphicspath{{C:/Users/david/coding/02_TRAFFIC_BIAS_paper/6000_LaTeX/figures/}}

\begin{document}

\title{Preliminary Gas Elasticity Heterogeneity Memo}
\maketitle



\section{Abstract}



\section{Introduction}
Following from Bates \& Kim (JAE 2024), this paper considers underlying characteristics that lead to the observed state heterogeneity in gas price elasticity in the second-stage of their Per-Cluster Intrumental Variables (PCIV) estimation. \\

To this end, we use the following state characteristics as predictors, preferring 1980 measures to avoid endogeneity concerns: \\



\section{Data Description}
% Describing Covariates
\begin{itemize}
	\item Minimum temperature in 1980. 
	\item Average temperature in 1980. 
	\item Maximum temperature in 1980. 
	\item Inches of precipitation in 1980. 
	\item Percentage flatness (Dobson \& Campbell 2014). 
	\item Miles of land area. 
	\item Meters of border with a different state. 
	\item Percentage of state population living in a state with a border county in 1980. 
	\item Percentage of state population utilizing public transport in 1980. 
\end{itemize}

\begin{center}
	\input gas_elasticity_dstats.tex
\end{center}



\section{Methodology \& Results}
%	Estimating Equation 
\begin{equation} 
	Elasticity_{s} = \boldsymbol{X} \boldsymbol{\beta} + u_{s} 
\end{equation}

%	Main Predictors, Weighted
\begin{table}[H]
	\begin{center}
		\caption{Main Predictors, Weighted}
		{
\def\sym#1{\ifmmode^{#1}\else\(^{#1}\)\fi}
\begin{tabular}{l*{6}{c}}
\hline\hline
            &\multicolumn{6}{c}{Gas Price Elasticity}                                                                                           \\
            &\multicolumn{1}{c}{(1)}         &\multicolumn{1}{c}{(2)}         &\multicolumn{1}{c}{(3)}         &\multicolumn{1}{c}{(4)}         &\multicolumn{1}{c}{(5)}         &\multicolumn{1}{c}{(6)}         \\
\hline
Land Area   &       0.838\sym{***}&                     &                     &                     &                     &       0.391\sym{**} \\
            &     (0.205)         &                     &                     &                     &                     &     (0.176)         \\
[1em]
Weather Ind.&                     &       0.050\sym{***}&                     &                     &                     &       0.036\sym{***}\\
            &                     &     (0.016)         &                     &                     &                     &     (0.013)         \\
[1em]
Urbanization&                     &                     &      -0.286\sym{***}&                     &                     &      -0.098         \\
            &                     &                     &     (0.106)         &                     &                     &     (0.115)         \\
[1em]
Pop. Density&                     &                     &                     &      -0.011\sym{***}&                     &      -0.008\sym{***}\\
            &                     &                     &                     &     (0.002)         &                     &     (0.002)         \\
[1em]
Cheaper Border&                     &                     &                     &                     &      -0.202\sym{**} &      -0.022         \\
            &                     &                     &                     &                     &     (0.094)         &     (0.084)         \\
\hline
Weights     &         Yes         &         Yes         &         Yes         &         Yes         &         Yes         &         Yes         \\
Locales:    &      51.000         &      51.000         &      51.000         &      51.000         &      51.000         &      51.000         \\
\hline\hline
\end{tabular}
}

%		\caption*{}
	\end{center}
\end{table}

%	Environmental Predictors, Weighted
\begin{table}[H]
	\begin{center}
		\caption{Environmental Predictors, Weighted}
		{
\def\sym#1{\ifmmode^{#1}\else\(^{#1}\)\fi}
\begin{tabular}{l*{3}{c}}
\hline\hline
            &\multicolumn{3}{c}{Gas Price Elasticity}                         \\
            &\multicolumn{1}{c}{(1)}         &\multicolumn{1}{c}{(2)}         &\multicolumn{1}{c}{(3)}         \\
\hline
Land Area   &       0.838\sym{***}&                     &                     \\
            &     (0.205)         &                     &                     \\
[1em]
Weather Ind.&                     &       0.050\sym{***}&                     \\
            &                     &     (0.016)         &                     \\
[1em]
Flatness Ind.&                     &                     &       0.016         \\
            &                     &                     &     (0.021)         \\
\hline
Weights     &         Yes         &         Yes         &         Yes         \\
Locales:    &      51.000         &      51.000         &      49.000         \\
\hline\hline
\end{tabular}
}

%		\caption*{}
	\end{center}
\end{table}

%	Transportation Predictors, Weighted
\begin{table}[H]
	\begin{center}
		\caption{Transportation Predictors, Weighted}
		{
\def\sym#1{\ifmmode^{#1}\else\(^{#1}\)\fi}
\begin{tabular}{l*{4}{c}}
\hline\hline
            &\multicolumn{4}{c}{Gas Price Elasticity}                                               \\
            &\multicolumn{1}{c}{(1)}         &\multicolumn{1}{c}{(2)}         &\multicolumn{1}{c}{(3)}         &\multicolumn{1}{c}{(4)}         \\
\hline
Urban Roads &       0.024         &                     &                     &                     \\
            &     (0.116)         &                     &                     &                     \\
[1em]
All Roads   &                     &       0.037         &                     &                     \\
            &                     &     (0.038)         &                     &                     \\
[1em]
Urbanization&                     &                     &      -0.286\sym{***}&                     \\
            &                     &                     &     (0.106)         &                     \\
[1em]
Pub. Trans. &                     &                     &                     &      -0.560         \\
            &                     &                     &                     &     (0.407)         \\
\hline
Weights     &         Yes         &         Yes         &         Yes         &         Yes         \\
Locales:    &      51.000         &      51.000         &      51.000         &      51.000         \\
\hline\hline
\end{tabular}
}

%		\caption*{}
	\end{center}
\end{table}

%	Demographic Predictors, Weighted
\begin{table}[H]
	\begin{center}
		\caption{Demographic Predictors, Weighted}
		{
\def\sym#1{\ifmmode^{#1}\else\(^{#1}\)\fi}
\begin{tabular}{l*{4}{c}}
\hline\hline
            &\multicolumn{4}{c}{Gas Price Elasticity}                                               \\
            &\multicolumn{1}{c}{(1)}         &\multicolumn{1}{c}{(2)}         &\multicolumn{1}{c}{(3)}         &\multicolumn{1}{c}{(4)}         \\
\hline
Population  &      -0.001         &                     &                     &                     \\
            &     (0.003)         &                     &                     &                     \\
[1em]
Pop. Density&                     &      -0.011\sym{***}&                     &                     \\
            &                     &     (0.002)         &                     &                     \\
[1em]
Border Pop. &                     &                     &      -0.168\sym{**} &                     \\
            &                     &                     &     (0.066)         &                     \\
[1em]
Licensure   &                     &                     &                     &       0.780\sym{*}  \\
            &                     &                     &                     &     (0.415)         \\
\hline
Weights     &         Yes         &         Yes         &         Yes         &         Yes         \\
Locales:    &      51.000         &      51.000         &      51.000         &      51.000         \\
\hline\hline
\end{tabular}
}

%		\caption*{}
	\end{center}
\end{table}

%%%%%%%%%%

%	Main Predictors, Unweighted
\begin{table}[H]
	\begin{center}
		\caption{Main Predictors, Unweighted}
		{
\def\sym#1{\ifmmode^{#1}\else\(^{#1}\)\fi}
\begin{tabular}{l*{6}{c}}
\hline\hline
            &\multicolumn{6}{c}{Gas Price Elasticity}                                                                                           \\
            &\multicolumn{1}{c}{(1)}         &\multicolumn{1}{c}{(2)}         &\multicolumn{1}{c}{(3)}         &\multicolumn{1}{c}{(4)}         &\multicolumn{1}{c}{(5)}         &\multicolumn{1}{c}{(6)}         \\
\hline
Land Area   &       0.442         &                     &                     &                     &                     &       0.203         \\
            &     (0.405)         &                     &                     &                     &                     &     (0.203)         \\
[1em]
Weather Ind.&                     &       0.016         &                     &                     &                     &       0.026\sym{*}  \\
            &                     &     (0.015)         &                     &                     &                     &     (0.015)         \\
[1em]
Urbanization&                     &                     &      -0.663\sym{***}&                     &                     &      -0.224\sym{*}  \\
            &                     &                     &     (0.241)         &                     &                     &     (0.122)         \\
[1em]
Pop. Density&                     &                     &                     &      -0.010\sym{***}&                     &      -0.008\sym{***}\\
            &                     &                     &                     &     (0.000)         &                     &     (0.001)         \\
[1em]
Cheaper Border&                     &                     &                     &                     &      -0.199         &       0.036         \\
            &                     &                     &                     &                     &     (0.163)         &     (0.073)         \\
\hline
Weights     &          No         &          No         &          No         &          No         &          No         &          No         \\
Locales:    &      51.000         &      51.000         &      51.000         &      51.000         &      51.000         &      51.000         \\
\hline\hline
\end{tabular}
}

%		\caption*{}
	\end{center}
\end{table}

%	Environmental Predictors, Unweighted
\begin{table}[H]
	\begin{center}
		\caption{Environmental Predictors, Unweighted}
		\input{etable_y1_environmental.tex}
%		\caption*{}
	\end{center}
\end{table}

%	Transportation Predictors, Unweighted
\begin{table}[H]
	\begin{center}
		\caption{Transportation Predictors, Unweighted}
		\input{etable_y1_transportation.tex}
%		\caption*{}
	\end{center}
\end{table}

%	Demographic Predictors, Unweighted
\begin{table}[H]
	\begin{center}
		\caption{Demographic Predictors, Unweighted}
		\input{etable_y1_demographic.tex}
%		\caption*{}
	\end{center}
\end{table}

\end{document}
